\documentclass[preview]{standalone}

\usepackage{amsmath}
\usepackage{amssymb}
\usepackage{stellar}
\usepackage{definitions}

\begin{document}

\id{cyclic-subspaces-companion-matrices}
\genpage

\section{Torsion and cyclic subspaces}

\begin{snippetproposition}{finite-dimensional-endomorphism-module-torsion}{The associated polynomial module is torsion}
    Let \(V\) be a finite-dimensional \(K\)-\vectorspace and let
    \(\varphi\colon V\fromto V\) be a \linearendomorphism. Then the
    \endomorphismmodule is torsion:
    \[
        V=\moduletor(V).
    \]
\end{snippetproposition}

\begin{snippetproof}{finite-dimensional-endomorphism-module-torsion-proof}{finite-dimensional-endomorphism-module-torsion}{The associated polynomial module is torsion}
    Let \(v\neq0_V\). Finite dimensionality makes the sequence
    \[
        v,\varphi(v),\varphi^2(v),\ldots
    \]
    linearly dependent. Let \(m\geq1\) be minimal such that
    \(v,\ldots,\varphi^m(v)\) are dependent. Then
    \(v,\ldots,\varphi^{m-1}(v)\) are independent, so a dependence
    relation must have a nonzero coefficient at \(\varphi^m(v)\).
    After normalization, it has the form
    \[
        a_0v+a_1\varphi(v)+\cdots+a_{m-1}\varphi^{m-1}(v)
        +\varphi^m(v)=0_V.
    \]
    Thus the nonzero monic polynomial
    \[
        f_v=a_0+a_1X+\cdots+a_{m-1}X^{m-1}+X^m
    \]
    lies in \(\moduleann(v)\). Hence every nonzero vector is torsion,
    and \(0_V\) is torsion as well.
\end{snippetproof}

\begin{snippetdefinition}{endomorphism-cyclic-subspace-definition}{Cyclic subspace}
    For \(v\in V\), the \emph{cyclic subspace generated by \(v\)} is
    \[
        W(v)=K[X]v
        =
        \{f\cdot v\suchthat f\in K[X]\}.
    \]
    It is the smallest \(\varphi\)-\invariantsubspace containing \(v\).
\end{snippetdefinition}

\begin{snippetproposition}{cyclic-subspace-annihilator-basis}{Annihilator and basis of a cyclic subspace}
    Let \(v\neq0_V\), and let \(m\) be minimal such that
    \(v,\varphi(v),\ldots,\varphi^m(v)\) are linearly dependent. There
    exists a unique monic polynomial
    \[
        f_v=X^m+a_{m-1}X^{m-1}+\cdots+a_1X+a_0
    \]
    of least degree satisfying \(f_v\cdot v=0_V\). Moreover,
    \[
        \moduleann(v)=(f_v)
    \]
    and
    \[
        W(v)
        =
        \operatorname{Span}_K
        \{v,\varphi(v),\ldots,\varphi^{m-1}(v)\}.
    \]
    The displayed vectors form a basis of \(W(v)\).
\end{snippetproposition}

\begin{snippetproof}{cyclic-subspace-annihilator-basis-proof}{cyclic-subspace-annihilator-basis}{Annihilator and basis of a cyclic subspace}
    The torsion proof constructs \(f_v\), and minimality makes the
    vectors through \(\varphi^{m-1}(v)\) independent. If
    \(g\in\moduleann(v)\), divide by \(f_v\):
    \[
        g=qf_v+r,
        \qquad
        r=0
        \quad\text{or}\quad
        \deg r<m.
    \]
    Acting on \(v\) gives \(r\cdot v=0_V\). A nonzero remainder would
    give a dependence among
    \(v,\varphi(v),\ldots,\varphi^{m-1}(v)\), so \(r=0\). Therefore
    \(f_v\divides g\) and \(\moduleann(v)=(f_v)\). The monic generator
    of a nonzero ideal of \(K[X]\) is unique.

    The relation \(f_v\cdot v=0_V\) expresses \(\varphi^m(v)\) as a
    combination of the preceding iterates. Applying \(\varphi\)
    repeatedly does the same for every later iterate, proving the span
    formula.
\end{snippetproof}

\begin{snippetcorollary}{cyclic-subspace-polynomial-quotient}{A cyclic subspace as a polynomial quotient}
    With the preceding notation,
    \[
        W(v)\moduleisomorphic K[X]/(f_v)
    \]
    as \(K[X]\)-modules.
\end{snippetcorollary}

\begin{snippetproof}{cyclic-subspace-polynomial-quotient-proof}{cyclic-subspace-polynomial-quotient}{A cyclic subspace as a polynomial quotient}
    The surjective \modulehomomorphism
    \[
        K[X]\fromto W(v),
        \qquad
        f\mapsto f\cdot v,
    \]
    has kernel \(\moduleann(v)=(f_v)\). Apply the first isomorphism
    theorem.
\end{snippetproof}

\section{Companion matrices}

\begin{snippetdefinition}{companion-matrix-definition}{Companion matrix}
    Let
    \[
        f=X^m+a_{m-1}X^{m-1}+\cdots+a_1X+a_0\in K[X]
    \]
    be monic of degree \(m\geq1\). Its \emph{companion matrix} is
    \[
        C(f)
        =
        \begin{pmatrix}
            0&0&\cdots&0&-a_0\\
            1&0&\cdots&0&-a_1\\
            0&1&\cdots&0&-a_2\\
            \vdots&\vdots&\ddots&\vdots&\vdots\\
            0&0&\cdots&1&-a_{m-1}
        \end{pmatrix}.
    \]
\end{snippetdefinition}

\begin{snippetproposition}{cyclic-subspace-companion-matrix}{Matrix of an endomorphism on a cyclic subspace}
    Let \(f_v\) be the monic generator of \(\moduleann(v)\), with
    \(\deg f_v=m\). In the basis
    \[
        \mathcal B_v
        =
        \{v,\varphi(v),\ldots,\varphi^{m-1}(v)\}
    \]
    of \(W(v)\), the matrix of
    \(\varphi|_{W(v)}\colon W(v)\fromto W(v)\) is
    \[
        C(f_v).
    \]
\end{snippetproposition}

\begin{snippetproof}{cyclic-subspace-companion-matrix-proof}{cyclic-subspace-companion-matrix}{Matrix of an endomorphism on a cyclic subspace}
    For \(0\leq i\leq m-2\),
    \(\varphi(\varphi^i(v))=\varphi^{i+1}(v)\), which gives the first
    \(m-1\) columns. The relation
    \[
        f_v\cdot v=0_V
    \]
    gives
    \[
        \varphi^m(v)
        =
        -a_0v-a_1\varphi(v)-\cdots-a_{m-1}\varphi^{m-1}(v),
    \]
    which is the final column of \(C(f_v)\).
\end{snippetproof}

\begin{snippetlemma}{companion-matrix-characteristic-polynomial}{Characteristic polynomial of a companion matrix}
    For every monic polynomial \(f\in K[X]\),
    \[
        \chi_{C(f)}(T)
        =
        \det(TI_m-C(f))
        =
        f(T).
    \]
\end{snippetlemma}

\begin{snippetproof}{companion-matrix-characteristic-polynomial-proof}{companion-matrix-characteristic-polynomial}{Characteristic polynomial of a companion matrix}
    The claim is immediate for \(m=1\). For \(m\geq2\),
    \[
        TI_m-C(f)
        =
        \begin{pmatrix}
            T&0&\cdots&0&a_0\\
            -1&T&\cdots&0&a_1\\
            0&-1&\ddots&0&a_2\\
            \vdots&\vdots&\ddots&T&\vdots\\
            0&0&\cdots&-1&T+a_{m-1}
        \end{pmatrix}.
    \]
    Expanding along the first row and applying induction to the first
    minor gives
    \[
        \det(TI_m-C(f))
        =
        T\left(T^{m-1}+a_{m-1}T^{m-2}+\cdots+a_1\right)+a_0,
    \]
    which is \(f(T)\).
\end{snippetproof}

\begin{snippetnote}{cyclic-decomposition-companion-blocks}{Cyclic decompositions and companion blocks}
    Every \cyclicsubspace supplies a basis in which the endomorphism is
    represented by a \companionmatrix. Thus a decomposition of the
    associated \(K[X]\)-module into cyclic torsion modules becomes a
    block-diagonal decomposition into companion matrices.
\end{snippetnote}

\end{document}
